\subsubsection*{2.1.7.}
现在回顾 $R\Psi$ 的函子性质, 即将 1.3.5--1.3.10 导出. 设 $f:X\longrightarrow X'$ 为 $S$-概形态射, 沿用 1.3.5 的记号.

态射 (1.3.6.1) 定义
\begin{equation}\tag{2.1.7.1}
R\Psi\,Rf_*\longrightarrow Rf_*\,R\Psi.
\end{equation}
由固有基变换定理（SGA~4, XII, 5.1）, \underline{若 $f$ 固有, 则此态射是同构.}

态射 (1.3.7.1) 定义
\begin{equation}\tag{2.1.7.2}
f^*R\Psi\longrightarrow R\Psi f^*.
\end{equation}
由光滑基变换定理（SGA~4, XVI, 1.2）, \underline{若 $f$ 光滑, 则此态射是同构.}

$f$ 固有时的同构 (2.1.7.1) 与态射 (1.3.8.1) 定义一个态射
\begin{equation}\tag{2.1.7.3}
Rf_!R\Psi\longrightarrow R\Psi Rf_!,
\end{equation}
当 $f$ 固有时, 它是 (2.1.7.1) 的逆. 原印本此处误指向后面的特例 (2.1.8.1).

当 $f$ 拟有限时, (1.3.9.1) 导出为
\begin{equation}\tag{2.1.7.4}
R\Psi Rf^!\longrightarrow Rf^!R\Psi,
\end{equation}
当 $f$ 平展时, 它是 (2.1.7.2) 的逆. 原印本指向 (2.1.7.1), 这与所涉及的两个函子不相容.

设 $f:S'\longrightarrow S$ 为离散赋值环谱的变换, 置 $X'=X\times_SS'$:
\begin{center}
\begin{tikzcd}
X' \arrow[r, "f"] \arrow[d, "p"'] & X \arrow[d, "p"] \\
S' \arrow[r, "f"] & S
\end{tikzcd}
\end{center}
（记号中的歧义见 1.3.5）. 假设 $f$ 满射. 基变换箭头 (1.3.10.1) 导出为
\begin{equation}\tag{2.1.7.5}
f^*R\Psi\longrightarrow R\Psi f^*.
\end{equation}
对 $R\Psi_\eta$ 有同样的陈述.

\subsubsection*{2.1.8.}
(2.1.7.1) 与 (2.1.7.3) 的关键情形为 $X'=S$. 对 $K\in\operatorname{Ob}D^+(X_{\bar\eta},\Lambda)$, 此时得到态射
\begin{equation}\tag{2.1.8.1}
Rf_*\longrightarrow Rf_*R\Psi
\end{equation}
及
\begin{equation}\tag{2.1.8.2}
Rf_!R\Psi\longrightarrow Rf_!,
\end{equation}
因而得到
\begin{equation}\tag{2.1.8.3}
\mathbb H^i(X_{\bar\eta},K)
\longrightarrow
\mathbb H^i(X_{\bar s},R\Psi_{\bar\eta}(K))
\end{equation}
及
\begin{equation}\tag{2.1.8.4}
\mathbb H_c^i(X_{\bar s},R\Psi_{\bar\eta}(K))
\longrightarrow
\mathbb H_c^i(X_{\bar\eta},K),
\end{equation}
以及下列图的交换性
\begin{equation}\tag{2.1.8.5}
\begin{tikzcd}[column sep=large]
(R^if_*K)_{\bar s} \arrow[r, "\mathrm{sp}"] \arrow[d]
  & \mathbb H^i(X_{\bar\eta},K) \arrow[d] \\
\mathbb H^i(X_{\bar s},K) \arrow[r]
  & \mathbb H^i(X_{\bar s},R\Psi_{\bar\eta}(K))
\end{tikzcd}
\end{equation}
和
\begin{equation}\tag{2.1.8.6}
\begin{tikzcd}[column sep=large]
\mathbb H_c^i(X_{\bar s},K) \arrow[rr, "\mathrm{sp}"] \arrow[dr]
  & {} & \mathbb H_c^i(X_{\bar\eta},K) \\
{} & \mathbb H_c^i(X_{\bar s},R\Psi_{\bar\eta}(K)) \arrow[ur] & {}
\end{tikzcd}
.\end{equation}

对 $K\in\operatorname{Ob}D^+(X_\eta,\Lambda)$ 与 $\sigma\in I$, 下列图交换:
\begin{equation}\tag{2.1.8.7}
\begin{tikzcd}[column sep=small, cells={nodes={font=\small}}]
\mathbb H^i(X_{\bar\eta},K)
  \arrow[rr, dash, "\sigma-1"] \arrow[d]
  & {} & \mathbb H^i(X_{\bar\eta},K) \arrow[d] \\
\mathbb H^i(X_{\bar s},R\Psi_{\bar\eta}(K)) \arrow[r]
  & \mathbb H^i(X_{\bar s},R\Phi_{\bar\eta}(K))
      \arrow[r, "\operatorname{Var}(\sigma)"]
  & \mathbb H^i(X_{\bar s},R\Psi_{\bar\eta}(K))
\end{tikzcd}
\end{equation}
和
\begin{equation}\tag{2.1.8.8}
\begin{tikzcd}[column sep=small, cells={nodes={font=\small}}]
\mathbb H_c^i(X_{\bar\eta},K)
  \arrow[rr, "\sigma-1"]
  & {} & \mathbb H_c^i(X_{\bar\eta},K) \\
\mathbb H_c^i(X_{\bar s},R\Psi_{\bar\eta}(K)) \arrow[r] \arrow[u]
  & \mathbb H_c^i(X_{\bar s},R\Phi_{\bar\eta}(K))
      \arrow[r, "\operatorname{Var}(\sigma)"]
  & \mathbb H_c^i(X_{\bar s},R\Psi_{\bar\eta}(K)) \arrow[u]
\end{tikzcd}
.\end{equation}
在 (2.1.8.8) 中, 英文层恢复了原印本左下角 $X_{\bar s}$ 漏掉的上划线, 以及下方中央项漏掉的右括号. 但在 (2.1.8.7) 中, 上方连线确实印作不带箭头的线段, 而 (2.1.8.8) 则有箭头; 这里保留了原文可见的这种区别.

当 $f$ 固有时, (2.1.8.1) 与 (2.1.8.2), 以及 (2.1.8.3) 与 (2.1.8.4), 都是互逆的同构. 区别三角 (2.1.2.4) 给出长正合列
\begin{equation}\tag{2.1.8.9}
\cdots\longrightarrow\mathbb H^i(X_{\bar s},K)
\longrightarrow\mathbb H^i(X_{\bar\eta},K)
\longrightarrow\mathbb H^i(X_{\bar s},R\Phi_{\bar\eta}(K))
\longrightarrow\cdots.
\end{equation}
因此消灭复形 $R\Phi(K)$ 表达了以 $K$ 的逆像为系数的几何纤维 $X_{\bar s}$ 与 $X_{\bar\eta}$ 的上同调之差异. 由 (2.1.8.7), 变差 $\operatorname{Var}(\sigma)$ 决定惯性群 $I$ 在 $\mathbb H^i(X_{\bar\eta},K)$ 上的作用. 这些事实是消灭圈方法的所有应用的基础.
